\documentclass[11pt,a4paper,roman]{moderncv}
\usepackage[english]{babel}
\usepackage{ragged2e}
\usepackage{float}
\usepackage{graphicx}
\usepackage[utf8]{inputenc}


\moderncvstyle{classic}
\moderncvcolor{green} % Bullet point color

% Page margins
\usepackage[scale=0.8]{geometry} % Page margins

% Your Information, please revise
\name{Bojie}{Fu}
\address{Haidian District, Beijing, China}{No.19, Xinjiekouwai St.}
\phone[mobile]{+86 139 1062 3747}
\email{bfu@rcees.ac.cn}


\begin{document}

% Insert Olin Logo
\begin{minipage}[t]{\textwidth}
\includegraphics[width=0.40\textwidth]{bnu}
\end{minipage}


\recipient{Dear Editor in Chief,}{}
\opening{\vspace*{-2em}}
\closing{Sincerely,}{\vspace*{-2em}}
\enclosure[Enclosures]{Main manuscript, Appendix}
\makelettertitle

I am pleased to submit an original research article entitled ``Identifying regime transitions for water governance at a basin scale'' for consideration for publication in \textit{People and Nature} as a \textit{Research Article}.

Water governance determines ``who gets water, when, and how'' in most large river basins; missing governance means missing sustainability. However, the lack of a comprehensive but straightforward approach to identifying the changes in water governance presents a challenge for efforts to underpin it. Therefore, we choose indicators for the corresponding aspects (water stress, water services purpose, and water allocation) and combine them into an integrated water governance index (IWGI) to analyze long-term changes in a large river basin.

Applying the IWGI to the rapidly-changing Yellow River Basin (YRB) in China clarifies shifts in water governance between massive supply, transformation governance, and adaptation-oriented regimes. After interpreting these regimes with underlying causes, we propose a potentially widespread transition schema in a human-dominated hydrosocial cycle. Our analysis approach links river basin governance with transitions of hydrosocial water cycles and offers valuable sustainability guidelines for big river basins worldwide. Therefore, we believe that this manuscript is appropriate for publication by \textit{People and Nature} because it tightly echoes the widespread reader interest in environmental governance.

This manuscript has not been published and is not under consideration for publication elsewhere, and all the authors have read and approved the manuscript. The manuscript includes the main body of text with $4176$ words and $4$ figures; the three-section appendix has $3$ supporting figures and $2$ supporting tables. Furthermore, we declare no conflicts of interest to disclose.

If you feel that the manuscript is appropriate for your journal, we suggest the following reviewers:

\begin{itemize}
\item Fuqiang Tian (Tsinghua University); Email: tianfq@tsinghua.edu.cn; Dr. Tian focuses on socio-hydrology and is familiar with the case basin in this study.
\item Tom Gleeson (University of Victoria); Email: tgleeson@uvic.ca; Dr. Gleeson cares about the significant changes during the transformation toward a human-interrupted world, especially the hydrology system.
\item Murugesu Sivapalan (University of Illinois, Urbana-Champaign); Email: sivapala@illinois.edu; Dr. Sivapalan focuses on socio-hydrology, linking water and human society with basin studies.
\item Olli Varis (Aalto University); Email: olli.varis@aalto.fi; Dr. Varis analyzed global water-based social-ecological vulnerability by integrated index, and he is also good at water resources research.
\end{itemize}

\vspace{0.3cm}

I appreciate your consideration!

\vspace{0.3cm}

Sincerely,

\textbf{Bojie Fu} (On behalf of the author team)

%\begin{minipage}[t]{\textwidth}
%	\includegraphics[height=1cm]{signature.jpeg}
%\end{minipage}

State Key Laboratory of Earth Surface Processes and Resource Ecology,

Faculty of Geographical Science,

Beijing Normal University,

Beijing 100875, China

Email: bfu@rcees.ac.cn

\end{document}
